\subsection{Simulation Framework}
\label{sub:sim_framework}

All of our numerical results are produced with the \texttt{sdim} qudit stabilizer simulator \cite{sdim}, version 1.4.0, which extends tableau-based stabilizer simulation to prime-dimensional qudits and samples noisy circuits in the Pauli frame picture. \autoref{sec:sdim-dem} reviews the simulator in more detail. The circuits in this work are expressed in terms of the generalized Hadamard, phase, SUM and controlled-phase gates of \autoref{sec:bg}, the Pauli and SWAP gates, and their inverses, computational-basis measurements, ancilla resets, and the one- and two-qutrit Pauli noise channels of \autoref{sub:noise_models}.

For each random sequence we draw the Cliffords for the largest depth once and build every depth from them: depth $m$ uses $\bar C_1, \dots, \bar C_m$ and the inverses $\bar C_m^{-1}, \dots, \bar C_1^{-1}$. For each sequence and depth we generate one annotated \texttt{sdim} circuit with check rounds, which serves every uniform-interval schedule, and one check-free circuit with the same Cliffords for terminal postselection. The physical RB circuits use their own random sequences. Each circuit contains the scheduled logical gates and check rounds, the noise operations of \autoref{sub:noise_models}, and \texttt{DETECTOR} and \texttt{LOGICAL\_OBSERVABLE} declarations. A detector is a linear combination of measurement outcomes that is deterministic in the absence of noise. Here each detector is one ancilla readout of one check round, or in terminal postselection one classically evaluated $X$-type stabilizer. The logical observable is the terminal readout of $\bar X$. Each noise operation stores its probability as a multiple of $p$ (\autoref{tab:circuit-noise}). For every value of $p$ we set each probability to that multiple of $p$, so one generated circuit structure per sequence and depth serves the entire sweep. The Clifford catalog is built with the Cayley graph construction of \autoref{sec:clifford-gen}, and the inverse of each element reverses its word and inverts every symbol, without resynthesis. The mirror inverse applies the looked-up inverse of each sampled Clifford, one Clifford at a time in reverse order, rather than a single recovery gate.

The annotated circuits can be sampled with two backends. The first is direct Pauli-frame simulation. The second compiles each circuit into a compact $q$-ary detector error model (DEM), a structure introduced for qubit stabilizer circuits by the Stim simulator \cite{Gidney2021Stim}. We use the native DEM of \texttt{sdim} 1.4.0, \texttt{sdim.dem.DetectorErrorModel}, which turns each noise operation that a detector or the logical observable can see into one error mechanism and represents the depolarizing and shift channels exactly. Postselected LRB-D statistics can then be sampled at a cost that no longer depends on the size of the underlying stabilizer computation. All four production runs, the one- and two-ancilla layouts under SI1000 and SD6, use the DEM backend for the LRB-D, terminal postselection and physical RB circuits, and so do the four high-$p$ extension runs, one per layout and noise model (\autoref{sub:results_setup}). %
% DEM versus frame sampler for the two-ancilla (conventional-site) code: paper-rework/scratch/writer-conv/demcheck/run_check.py (run_check.out); one-ancilla code: tests/test_dem_backend.py
On depth-$2$ circuits of both layouts under both models at $p = 0.01$, the per-round pass rates and logical outcomes from the two backends agree within five standard errors over $4 \times 10^4$ shots. \autoref{sec:sdim-dem} explains the construction, why it is exact for our noise models, and how it is sampled.
